\chapter{Chronic Insomnia Disorder}
\index{Chronic Insomnia Disorder}
\label{sec:Chronic Insomnia Disorder}

\section{Overview}
Chronic insomnia disorder is a sleep disorder defined by persistent difficulty initiating sleep, maintaining sleep, or early morning awakening with inability to return to sleep, despite adequate opportunity for sleep, occurring at least three nights per week for at least three months, with significant daytime impairment (fatigue, mood disturbance, thinking and memory dysfunction, or reduced performance). The prevalence is 6--10\% in the general population, with females affected 1.5--2 times more frequently than males. This condition is among the most common mental health disorders worldwide. It affects people across all age groups, socioeconomic backgrounds, and geographic regions. Mental health conditions affect thoughts, emotions, behaviors, and overall well-being. These conditions are among the most common health concerns worldwide, affecting people across all age groups. Mental health disorders result from complex interactions of biological, psychological, and social factors. Effective treatments are available, and recovery is possible with appropriate support. This condition affects a significant number of people worldwide and can range from mild to severe in presentation. Early recognition and appropriate management are essential for optimal outcomes. A thorough understanding of the underlying mechanisms guides treatment decisions. Patient education and self-management play important roles in long-term care.
\section{Causes}
This condition happens because of hyperarousal of the hypothalamic-pituitary-adrenal axis (elevated cortisol, ACTH), increased high-frequency EEG activity during sleep, activation of the sympathetic nervous system, and maladaptive conditioning (negative associations with the sleep environment). The underlying mechanisms involve complex interactions between genetic predisposition and environmental factors that disrupt normal physiological processes. Neurotransmitter imbalances (serotonin, dopamine, norepinephrine) play key roles in many mental health conditions. Genetic factors contribute to vulnerability, with heritability estimates varying by condition. Early life experiences, trauma, and chronic stress can alter brain development and function. Environmental factors including social support, socioeconomic status, and life events influence mental health. The underlying mechanisms involve complex interactions between genetic predisposition and environmental factors. Disruption of normal physiological processes leads to the characteristic features of the condition. Multiple contributing factors may act synergistically to trigger or worsen the condition.
\section{Risk Factors}
Risk factors include advanced age, female sex, shift work, medical and psychiatric comorbidities, and genetic predisposition.

\begin{itemize}[noitemsep]
  \item Family history of mental illness
  \item Trauma or adverse childhood experiences
  \item Chronic stress
  \item Substance use
  \item Social isolation
  \item Underlying medical conditions
  \item Genetic predisposition and family history
  \item Advancing age
  \item Lifestyle factors (diet, exercise, smoking, alcohol)
  \item Environmental exposures
  \item Medication use
\end{itemize}
\section{Signs and Symptoms}
You may experience difficulty falling asleep, frequent awakenings, early morning awakening, and daytime impairment.

\begin{itemize}[noitemsep]
  \item Pain or discomfort in the affected area
  \item Functional impairment affecting daily activities
  \item Changes in appearance or function of affected tissues
  \item Systemic symptoms may include fatigue, fever, or weight changes
  \item Symptoms may develop gradually or appear suddenly
\end{itemize}
\section{Progression and Stages}
The condition may progress through different stages of severity over time. Early stages often involve mild or subtle symptoms that may be overlooked. Moderate stages involve more noticeable symptoms affecting daily function. Advanced stages may involve significant impairment and complications.
\section{Complications}
\begin{itemize}[noitemsep]
  \item Disease progression leading to worsening symptoms
  \item Secondary infections or injuries
  \item Side effects of treatment
  \item Reduced quality of life
  \item Emotional and psychological impact
\end{itemize}
\section{Diagnosis}
\begin{itemize}[noitemsep]
  \item Doctors usually diagnose this based on your symptoms and a physical exam based on ICSD-3 and DSM-5 criteria, supplemented by sleep diaries and actigraphy
  \item Polysomnography is not routinely indicated but may evaluate comorbid sleep disorders.
  \item Comprehensive medical history and physical examination
  \item Laboratory tests to assess organ function
  \item Imaging studies to visualize affected structures
  \item Specialized testing depending on the affected system
  \item Consultation with relevant specialists
\end{itemize}
\section{Prevention}
\begin{itemize}[noitemsep]
  \item Maintain a healthy lifestyle with balanced diet and exercise
  \item Avoid known risk factors
  \item Attend regular medical check-ups for early detection
  \item Follow recommended screening guidelines
\end{itemize}
\section{Treatment Options}
Treatment focuses on thinking and memory behavioral therapy for insomnia (CBT-I) as first-line therapy---stimulus control, sleep restriction, thinking and memory restructuring, and sleep hygiene education---with pharmacotherapy (melatonin, ramelteon, orexin receptor antagonists, benzodiazepine receptor agonists, trazodone, doxepin) as adjunctive treatment. Psychotherapy (cognitive behavioral therapy, interpersonal therapy, psychodynamic therapy) Pharmacotherapy (antidepressants, antipsychotics, mood stabilizers, anxiolytics) Combined treatment approaches for optimal outcomes Hospitalization for acute crisis management Support groups and peer support programs Lifestyle interventions (exercise, sleep hygiene, stress management) Mindfulness and meditation practices Family therapy and psychoeducation Vocational and social rehabilitation Long-term maintenance therapy for chronic conditions

\begin{itemize}[noitemsep]
  \item Medications to manage symptoms and address underlying mechanisms
  \item Lifestyle modifications and self-care measures
  \item Physical therapy and rehabilitation
  \item Surgical intervention when indicated
  \item Regular monitoring and follow-up care
\end{itemize}
\section{Living with It}
\begin{itemize}[noitemsep]
  \item Follow your treatment plan as prescribed
  \item Maintain regular follow-up with your healthcare provider
  \item Make necessary lifestyle adjustments
  \item Seek support from family, friends, or support groups
  \item Stay informed about your condition
\end{itemize}
\section{Outlook}
\begin{itemize}[noitemsep]
  \item The outlook is generally positive with CBT-I
  \item Long-term outcomes favor CBT-I over medication alone.
\end{itemize}
